\section*{\centering{RESUMO}} 
A necessidade de odometria redundante com as mais diversas fontes de captação, se beneficia das recentes possibilidades de processamento rápido e eficiente da imagem em tempo real, por meio de algoritmos de classificação e regressão linear. O presente trabalho propõe uma abordagem baseada em Redes Neurais Convolutivas (RNC), utilizadas para detectar numericamente a distância entre o sensores de imagem e a superfície sobrevoada durante a operação de VANTs sobre a água. O sistema de RNC foi treinado com imagens sintéticas geradas em Blender, e avaliaram imagens reais capturadas por um quadrotor controlado remotamente, coletando, também, os dados de odometria inerentes e sincronizados para respectiva comparação. Os testes preliminares demonstraram a capacidade do sistema de reconhecer diferentes distâncias da superfície da água, mesmo com o conjunto de imagens reais em condição não ótima, em ambiente não controlado, com sujeiras e sombras, após o devido tratamento e filtragem descritos no escopo do trabalho.

{\bf Palavras-chave: VANT, offshore, aprendizagem de máquina, estimação de altitude} 


